\documentclass[a4paper, final]{article}
\usepackage[14pt]{extsizes} % для того чтобы задать нестандартный 14-ый размер шрифта
\usepackage{iftex}
\ifPDFTeX
\usepackage[T2A]{fontenc}
\usepackage[utf8]{inputenc}
\usepackage[russian]{babel}
\else
\usepackage{fontspec}
\IfFontExistsTF{Liberation Serif}{\setmainfont{Liberation Serif}}{\setmainfont{FreeSerif}}
\IfFontExistsTF{Liberation Sans}{\setsansfont{Liberation Sans}}{\setsansfont{FreeSans}}
\IfFontExistsTF{Liberation Mono}{\setmonofont{Liberation Mono}}{\setmonofont{FreeMono}}
\usepackage{babel}
\babelprovide[main,import]{russian}
\fi


\usepackage{ragged2e}
\usepackage{pdfpages}
\usepackage{tabularx}
\usepackage{array}
\usepackage{array,longtable,booktabs}
\usepackage{amsmath}
\usepackage{multirow}
\usepackage{multicol}
\usepackage{stackrel}
\usepackage[left=25mm, top=20mm, right=20mm, bottom=20mm, footskip=10mm]{geometry}
\usepackage{ragged2e} %для растягивания по ширине
\usepackage{sectsty}
\allsectionsfont{\raggedright}
\usepackage{setspace} %для межстрочного интервала
\usepackage{moreverb} %для работы с листингами
\usepackage{indentfirst} % для абзацного отступа
\usepackage{moreverb} %для печати в листинге исходного кода программ
\renewcommand\verbatimtabsize{4\relax}
\renewcommand\listingoffset{0.2em}
\renewcommand{\arraystretch}{1.4} % изменяю высоту строки в таблице
\usepackage[font=small, singlelinecheck=false, justification=raggedleft, format=plain, labelsep=period]{caption} %для настройки заголовка таблицы
\usepackage{graphicx} % Для вставки графических изображений
\graphicspath{{noiseimages/}}
\usepackage{color} %% это для отображения цвета в коде
\usepackage{xcolor} % цвета
\usepackage{listingsutf8}
\usepackage{hyperref}% для гиперссылок
\usepackage{enumitem} %для перечислений
\usepackage{pdflscape} %для pdf
\usepackage{pdfpages} %для pdf

\definecolor{apricot}{HTML}{FFF0DA}
\definecolor{mygreen}{rgb}{0,0.6,0}
\definecolor{string}{HTML}{B40000} % цвет строк в коде
\definecolor{comment}{HTML}{008000} % цвет комментариев в коде
\definecolor{keyword}{HTML}{1A00FF} % цвет ключевых слов в коде
\definecolor{morecomment}{HTML}{8000FF} % цвет include и других элементов в коде
\definecolor{captiontext}{HTML}{FFFFFF} % цвет текста заголовка в коде
\definecolor{captionbk}{HTML}{999999} % цвет фона заголовка в коде
\definecolor{bk}{HTML}{FFFFFF} % цвет фона в коде
\definecolor{frame}{HTML}{999999} % цвет рамки в коде
\definecolor{brackets}{HTML}{B40000} % цвет скобок в коде

\setlist[enumerate,itemize]{leftmargin=1.2cm}

\definecolor{linkblue}{RGB}{10,90,200}
\hypersetup{colorlinks,
	pdftitle={Оптимизация запросов в реляционных базах данных},
	pdfauthor={П. М. Лобанов},
	allcolors=linkblue}
% подгружаемые языки — подробнее в документации listings
\lstloadlanguages{Java}
% включаем кириллицу и добавляем кое−какие опции
%\lstset{language =[LaTeX] TeX, % выбираем язык по умолчанию
	%extendedchars=true , % включаем не латиницу
	%escapechar = | , % |«выпадаем» в LATEX|
	%frame=tb , % рамка сверху и снизу
	%commentstyle=\itshape , % шрифт для комментариев
	%stringstyle =\bfseries} % шрифт для строк

\textheight=24cm % высота текста
\textwidth=16cm % ширина текста
\oddsidemargin=0pt % отступ от левого края
\topmargin=-1.5cm % отступ от верхнего края
\parindent=24pt % абзацный отступ
\parskip=0pt % интервал между абзацами
\setlength{\emergencystretch}{4em}
\tolerance=2000 % терпимость к "жидким" строкам
\raggedbottom

\addto\captionsrussian{\def\refname{\hspace{1.15cm} Список литературы}}

\begin{document} % начало документа
	
	\lstset{
		language=Java, % Язык кода по умолчанию
		morekeywords={*,...}, % если хотите добавить ключевые слова, то добавляйте
		% Цвета
		keywordstyle=\color{keyword}\ttfamily\bfseries,
		%stringstyle=\color{string}\ttfamily,
		stringstyle=\ttfamily\color{red!50!brown},
		commentstyle=\color{comment}\ttfamily,
		morecomment=[l][\color{morecomment}]{\#},
		% Настройки отображения
		breaklines=true, % Перенос длинных строк
		basicstyle=\ttfamily\footnotesize, % Шрифт для отображения кода
		backgroundcolor=\color{bk}, % Цвет фона кода
		frame=lrb,xleftmargin=\fboxsep,xrightmargin=-\fboxsep, % Рамка, подогнанная к заголовку
		rulecolor=\color{frame}, % Цвет рамки
		tabsize=3, % Размер табуляции в пробелах
		% Настройка отображения номеров строк. Если не нужно, то удалите весь блок
		numbers=left, % Слева отображаются номера строк
		stepnumber=1, % Каждую строку нумеровать
		numbersep=5pt, % Отступ от кода
		numberstyle=\small\color{black}, % Стиль написания номеров строк
		% Для отображения русского языка
		extendedchars=true,
		literate={Ö}{ {\"O} }1
		{~}{ {\textasciitilde} }1
		{а}{ {\selectfont\char224} }1
		{б}{ {\selectfont\char225} }1
		{в}{ {\selectfont\char226} }1
		{г}{ {\selectfont\char227} }1
		{д}{ {\selectfont\char228} }1
		{е}{ {\selectfont\char229} }1
		{ё}{ {\"e} }1
		{ж}{ {\selectfont\char230} }1
		{з}{ {\selectfont\char231} }1
		{и}{ {\selectfont\char232} }1
		{й}{ {\selectfont\char233} }1
		{к}{ {\selectfont\char234} }1
		{л}{ {\selectfont\char235} }1
		{м}{ {\selectfont\char236} }1
		{н}{ {\selectfont\char237} }1
		{о}{ {\selectfont\char238} }1
		{п}{ {\selectfont\char239} }1
		{р}{ {\selectfont\char240} }1
		{с}{ {\selectfont\char241} }1
		{т}{ {\selectfont\char242} }1
		{у}{ {\selectfont\char243} }1
		{ф}{ {\selectfont\char244} }1
		{х}{ {\selectfont\char245} }1
		{ц}{ {\selectfont\char246} }1
		{ч}{ {\selectfont\char247} }1
		{ш}{ {\selectfont\char248} }1
		{щ}{ {\selectfont\char249} }1
		{ъ}{ {\selectfont\char250} }1
		{ы}{ {\selectfont\char251} }1
		{ь}{ {\selectfont\char252} }1
		{э}{ {\selectfont\char253} }1
		{ю}{ {\selectfont\char254} }1
		{я}{ {\selectfont\char255} }1
		{А}{ {\selectfont\char192} }1
		{Б}{ {\selectfont\char193} }1
		{В}{ {\selectfont\char194} }1
		{Г}{ {\selectfont\char195} }1
		{Д}{ {\selectfont\char196} }1
		{Е}{ {\selectfont\char197} }1
		{Ё}{ {\"E} }1
		{Ж}{ {\selectfont\char198} }1
		{З}{ {\selectfont\char199} }1
		{И}{ {\selectfont\char200} }1
		{Й}{ {\selectfont\char201} }1
		{К}{ {\selectfont\char202} }1
		{Л}{ {\selectfont\char203} }1
		{М}{ {\selectfont\char204} }1
		{Н}{ {\selectfont\char205} }1
		{О}{ {\selectfont\char206} }1
		{П}{ {\selectfont\char207} }1
		{Р}{ {\selectfont\char208} }1
		{С}{ {\selectfont\char209} }1
		{Т}{ {\selectfont\char210} }1
		{У}{ {\selectfont\char211} }1
		{Ф}{ {\selectfont\char212} }1
		{Х}{ {\selectfont\char213} }1
		{Ц}{ {\selectfont\char214} }1
		{Ч}{ {\selectfont\char215} }1
		{Ш}{ {\selectfont\char216} }1
		{Щ}{ {\selectfont\char217} }1
		{Ъ}{ {\selectfont\char218} }1
		{Ы}{ {\selectfont\char219} }1
		{Ь}{ {\selectfont\char220} }1
		{Э}{ {\selectfont\char221} }1
		{Ю}{ {\selectfont\char222} }1
		{Я}{ {\selectfont\char223} }1
		{\{}{ { {\color{brackets}\{} } }1 % Цвет скобок {
			{\} }{ { {\color{brackets}\} } } }1 % Цвет скобок }
	}
	
	
	% НАЧАЛО ТИТУЛЬНОГО ЛИСТА
	\thispagestyle{empty}
	\begin{center}
		\normalsize{МИНИСТЕРСТВО НАУКИ И ВЫСШЕГО ОБРАЗОВАНИЯ РОССИЙСКОЙ ФЕДЕРАЦИИ \\ ФЕДЕРАЛЬНОЕ ГОСУДАРСТВЕННОЕ АВТОНОМНОЕ ОБРАЗОВАТЕЛЬНОЕ УЧРЕЖДЕНИЕ ВЫСШЕГО \\ ОБРАЗОВАНИЯ\\ «Санкт-Петербургский политехнический университет Петра Великого»}
		\normalsize{Институт компьютерных наук и кибербезопасности}\\[3pt]
		\normalsize{Высшая школа технологий искусственного интеллекта}\\[10pt]
		\normalsize{Направление: 02.04.01 Математика и компьютерные науки}\\[10pt]
		\hfill \break
		\hfill \break
		\hfill \break
		\hfill \break
		\hfill \break
		\hfill \break
		\normalsize{Обзор литературы по теме:}\\
		\normalsize{«Оптимизация запросов в реляционных базах данных»}\\
		\normalsize{по дисциплине «Цифровые ресурсы в научном исследовании»}\\[10pt]
	\end{center}
	\hfill \break
	\hfill \break
	\hfill \break
	\hfill \break
	\noindent\begin{tabular*}{\linewidth}{@{}l@{\extracolsep{\fill}}ccl@{}}	
		\!\!\!Студент, & \hspace{0.6cm} & & \\
		\!\!\!группы 5140201/60301 & \hspace{0.6cm} & \underline{\hspace{3cm}} &Лобанов П. М. \\\\
		\!\!\!Преподаватель & \hspace{0.6cm}  & \underline{\hspace{3cm}} &\hspace{0pt} Моторин Д. Е. \\\\	
		&&
		
	\end{tabular*} 	
	\begin{flushright}
		<<\underline{\hspace{1cm}}>>\underline{\hspace{2.5cm}} 2026 г.
	\end{flushright}
	\hfill \break
	\hfill \break
	\hfill \break
	\begin{center} \small{Санкт-Петербург, 2026} \end{center}
	\newpage %Содержание% выключаем отображение номера для этой страницы
	\tableofcontents
	
	
	\newpage
	\section*{Введение}
	\addcontentsline{toc}{section}{Введение}
	
	При построении плана SQL-запроса оптимизатору приходится оценивать объём промежуточных результатов до выполнения соединений. Если после первого соединения остаётся значительно больше строк, чем ожидалось, выбранный способ обработки этих строк на следующем этапе может оказаться слишком затратным. На запросах к базе данных IMDb Leis и соавторы показали, что время выполнения зависит как от точности таких оценок, так и от того, какие порядки соединения рассмотрел оптимизатор~\cite{leis2015}.
	
	Классические и обучаемые методы решают эту задачу на разных этапах построения плана. Стоимостная модель позволяет сравнить варианты по ожидаемым затратам, а эвристический поиск --- ограничить число проверяемых порядков соединения. Обучаемые методы могут уточнять прогноз числа строк, направлять поиск с учётом результатов предыдущих запусков или предлагать изменения в уже построенном плане. Отсюда различие в оценке результатов: для модели кардинальности существенна ошибка прогноза, а для метода выбора плана --- время выполнения запроса~\cite{manegold2000,goncalves2014,wang2024,autoquo2024,foss2024}.
	
	При распределённой обработке даже для заданного порядка соединения остаётся выбрать, какие данные передавать между узлами и как распределить вычисления. Например, при перераспределении строк по ключу все строки с одним значением направляются одному исполнителю. Если это значение встречается особенно часто, остальные исполнители могут закончить работу раньше и ожидать завершения одной задачи. Поэтому для выбора способа соединения наряду с размерами таблиц могут потребоваться сведения о размещении строк и частотах отдельных ключей~\cite{liang2024,prpq2017,flowjoin2026}.
	
	Цель обзора --- сопоставить классические и обучаемые методы оптимизации SQL-запросов с соединениями и определить, как меняется выбор плана при переходе к распределённой обработке. В обзор включены 16 экспериментальных работ. Для каждой рассматриваются решение, которое принимает алгоритм, сведения, на которые он опирается, и показатель, по которому авторы оценивают результат. Отдельно рассматривается вопрос для дальнейшего исследования: при каких условиях учёт частот ключей ускоряет распределённое соединение настолько, чтобы выигрыш покрыл затраты на сбор статистики?
	
	\newpage
	\section{Основные задачи оптимизатора запросов}
	
	Один и тот же SQL-запрос может быть выполнен несколькими способами. Например, если запрос содержит несколько внутренних соединений, СУБД может изменить порядок соединения таблиц или использовать разные алгоритмы поиска совпадающих строк. Планы, которые для одних и тех же исходных данных возвращают предусмотренный запросом результат, называют семантически эквивалентными. При этом время их выполнения может существенно различаться. Классический стоимостный оптимизатор выбирает среди рассмотренных планов вариант с наименьшей расчётной стоимостью. При этом он оценивает размеры промежуточных результатов, рассчитывает стоимость операций и определяет, какие варианты плана следует проверить~\cite{leis2015}.
	
	В работе \textit{How Good Are Query Optimizers, Really?} Leis и соавторы исследовали, на каком из этих этапов возникают ошибки, из-за которых оптимизатор выбирает медленный план. Для экспериментов авторы разработали Join Order Benchmark (JOB) --- набор запросов с несколькими соединениями к базе данных фильмов IMDb. Чтобы отделить влияние ошибок в оценке числа строк от влияния стоимостной модели и алгоритма поиска, исследователи, в частности, подставляли фактическое число строк вместо прогнозируемого. Затем они проверяли, изменился ли выбранный план и как это сказалось на времени выполнения запроса~\cite{leis2015}.
	
	Первая задача --- оценить кардинальность, то есть число строк, которое останется после фильтрации или получится в результате соединения. Для этого оптимизатор использует статистику о данных. Например, по распределению значений в столбце можно оценить долю строк, удовлетворяющих отдельному условию. Если два условия связаны оператором \texttt{AND}, то при допущении об их независимости оптимизатор может оценить долю строк, удовлетворяющих обоим условиям, как произведение соответствующих долей. Если же значения столбцов взаимосвязаны, такая оценка может существенно отличаться от фактического результата. Отклонение повлияет и на оценку объёма данных, поступающих на вход следующей операции~\cite{leis2015,wang2024}.
	
	Вторая задача --- использовать ожидаемое число строк для оценки затрат на выполнение операции. Способ расчёта зависит от алгоритма соединения. Например, при соединении вложенными циклами с поиском по индексу для каждой строки внешней таблицы выполняется поиск совпадений во внутренней таблице. Если после фильтрации ожидается мало внешних строк, стоимостная модель учтёт лишь небольшое число таких обращений. При хеш-соединении затраты складываются иначе: по строкам одной таблицы строится хеш-таблица, а для строк другой в ней выполняется поиск совпадений.
	
	В экспериментах на JOB это различие проявилось при недооценке числа строк. Соединение вложенными циклами могло получить низкую расчётную стоимость, хотя при выполнении запроса обращений к индексу требовалось значительно больше, чем предполагал оптимизатор. Планы с хеш-соединениями и полным сканированием таблиц в тех же условиях оказались устойчивее к таким ошибкам. В целом на исследованной нагрузке, когда данные помещались в оперативной памяти, точность оценки числа строк влияла на время выполнения выбранных планов сильнее, чем различия между рассмотренными стоимостными моделями~\cite{leis2015}.
	
	Третья задача --- определить, какие планы сравнивать. С увеличением числа таблиц быстро растёт число допустимых порядков их соединения. Динамическое программирование позволяет повторно использовать планы, уже найденные для подмножеств таблиц, при построении плана всего запроса. Эвристические методы ограничивают число проверяемых вариантов и тем самым сокращают время поиска, но могут пропустить выгодный порядок соединения. В экспериментах Leis и соавторов более полный поиск позволял находить планы с меньшим фактическим временем выполнения даже при неточных оценках кардинальности~\cite{leis2015}. Поэтому при сравнении методов оптимизации важно учитывать как способ расчёта стоимости, так и то, какие варианты алгоритм включает в поиск.
	
	
	\newpage
	\section{Стоимостные и эвристические методы}
	
	Оценка кардинальности даёт исходные данные для сравнения планов. Даже если число строк известно точно, оптимизатору предстоит решить, какие порядки соединения проверить за отведённое время и как оценить затраты на каждый вариант. Первую задачу Gonçalves и соавторы изучали на примере эвристического поиска в СУБД H2. Вторую можно проследить на примере хеш-соединения в оперативной памяти и построения индекса для временной таблицы~\cite{goncalves2014,manegold2000,rahman2026}.
	
	\subsection{Поиск порядка соединения и оценка обращений к памяти}
	
	Gonçalves и соавторы исследовали, можно ли найти порядок соединения с минимальной расчётной стоимостью, проверив лишь часть допустимых вариантов. Для этого они разработали эволюционный алгоритм MAQO и реализовали его в оптимизаторе СУБД H2. Агент MAQO хранит один вариант порядка соединения таблиц и действует в соответствии с заданным профилем. Профиль определяет, какие операции агент применяет к решению: наряду со скрещиванием и мутацией в MAQO используются эвристики построения порядка и локальный поиск для его улучшения~\cite{goncalves2014}.
	
	Авторы отдельно проверяли способность MAQO находить глобальный минимум заданной функции стоимости и измеряли время выполнения найденных планов. В одном из приведённых примеров план MAQO имел меньшую расчётную стоимость и выполнился почти на 23\,\% быстрее плана штатного оптимизатора H2. Это результат отдельного запроса. Чтобы судить о качестве поиска на всей нагрузке, нужно дополнительно учитывать, как часто эвристика находит выгодный порядок и сколько времени тратит на его поиск~\cite{goncalves2014}.
	
	Почему два плана, обрабатывающих одинаковое число строк, могут иметь разную стоимость, видно из работы Manegold, Boncz и Kersten. Авторы исследовали соединение по равенству в СУБД Monet. Если хеш-таблица велика, поиск очередного ключа может потребовать обращения к участку памяти, которого сейчас нет в кеше процессора. Один из способов уменьшить такие затраты --- предварительно разбить входные данные на части по нескольким битам хеш-значения. Тогда соответствующие части двух таблиц можно соединять по очереди, работая с меньшей хеш-таблицей~\cite{manegold2000}.
	
	У такого разбиения есть собственная цена: строки нужно прочитать и распределить по частям. Чтобы оценить, когда эти затраты окупаются, авторы включили в модель стоимость обращений к разным уровням памяти. Прогнозы они проверяли в Monet по времени выполнения и показаниям аппаратных счётчиков, включая число промахов кешей первого и второго уровней. В совокупности с экспериментом MAQO это объясняет, почему поиск и стоимостную модель проверяют раздельно: эвристика определяет, какие планы будут рассмотрены, а модель --- какой из них получит преимущество при сравнении~\cite{manegold2000,goncalves2014}.
	
	\subsection{Выбор индекса для промежуточного результата}
	
	До сих пор речь шла об операциях внутри одного запроса. Rahman и соавторы рассмотрели случай, когда один оператор SQL-процедуры сохраняет промежуточный результат во временной таблице, а последующие операторы несколько раз читают из неё строки по одному условию. Индекс может ускорить каждое чтение, но его создание также требует времени. Поэтому авторы оценивали суммарные затраты за всё время использования таблицы~\cite{rahman2026}.
	
	Разработанная авторами система AutoCox определяет, какой оператор создаёт временную таблицу и какие операторы читают её впоследствии. Затем она сопоставляет два варианта. В первом временную таблицу читают без индекса. Во втором учитываются стоимость построения индекса и ожидаемые затраты на чтение по нему. Выбор зависит от числа строк, доли записей, удовлетворяющих условию, и количества повторных обращений~\cite{rahman2026}.
	
	В экспериментах авторы сообщают о снижении затрат на зависимые запросы на 67\,\% и об ускорении выполнения программы в 261 раз на одной из исследованных нагрузок. Первый показатель характеризует обращения к промежуточному результату, второй --- выполнение программы в целом. Затраты на построение индекса возникают один раз за выполнение процедуры, а экономия времени зависит от того, какие строки и сколько раз затем потребуются другим операторам. Поэтому выигрыш для программы с повторными чтениями нельзя переносить на произвольный одиночный запрос~\cite{rahman2026}.
	
	Следовательно, кардинальность задаёт объём работы. При сравнении вариантов оптимизатору также важны способ доступа к строкам и затраты на подготовку разбиения или индекса~\cite{manegold2000,rahman2026}.
	
	\newpage
	\section{Обучаемые методы оптимизации}
	
	В рассмотренных выше методах оптимизатор опирался на заданную функцию стоимости и правила поиска. В обучаемых методах часть решений зависит от закономерностей, которые модель выявляет на обучающих примерах. В одних работах модель предсказывает объём данных до выбора плана, а в других использует сведения о выполненных запросах, чтобы рекомендовать порядок таблиц или изменить уже выбранное соединение~\cite{wang2024,autoquo2024,foss2024}.
	
	\subsection{Оценка кардинальности запросов}
	
	Модель RvCE, предложенная Wang и соавторами, прогнозирует число строк в результате запроса. Для этого модель учитывает состав таблиц, связи между ними и структуру условий отбора. Например, условия, объединённые операторами \texttt{AND} и \texttt{OR}, по-разному определяют состав результата, даже если в обоих случаях сравниваются одни и те же столбцы. Поэтому RvCE обрабатывает условие как дерево логических операций, а связи между соединяемыми таблицами представляет в виде графа. Дополнительные сведения о данных модель получает из выборок~\cite{wang2024}.
	
	Авторы оценивают точность прогноза, в частности, по показателю $q$-error. Для положительных значений он равен отношению большего из двух чисел, прогнозируемого и фактического, к меньшему. Значение, равное единице, соответствует точному прогнозу. При этом завышение и занижение числа строк в одинаковое число раз дают одинаковый $q$-error, хотя последствия для плана могут различаться. В частности, недооценка числа строк может сделать привлекательным индексное соединение, которое при выполнении потребует большого числа обращений к индексу. На данных IMDb авторы получили более точные оценки, чем другие исследованные методы, обученные на запросах. Этот результат характеризует точность прогноза числа строк. Чтобы установить влияние на запрос в целом, требуется дополнительно проверить, изменились ли выбранный план и фактическое время его выполнения~\cite{wang2024,leis2015}.
	
	\subsection{Выбор и изменение плана по результатам выполнения}
	
	В AutoQuo результаты предыдущих запусков используются уже при выборе порядка соединения. Алгоритм строит частичный порядок, на каждом шаге рассматривая, какую таблицу добавить следующей. Поиск по дереву Монте-Карло позволяет исследовать новые продолжения и возвращаться к вариантам, которые ранее получили высокую оценку. Найденный частичный порядок передаётся PostgreSQL в виде подсказки, после чего штатный оптимизатор достраивает план. Обучаемая модель прогнозирует производительность кандидатов, а сведения о фактическом выполнении используются для её обновления. В экспериментах авторов на запросах к IMDb суммарное время выполнения сократилось на 29,73\,\% по сравнению с PostgreSQL~\cite{autoquo2024}.
	
	Для обучаемого поиска существенна стоимость получения опыта. В начале работы наблюдений ещё мало, а проверка неудачного плана может занять значительное время. В LOOPLINE для начального обучения используются сведения штатного оптимизатора PostgreSQL. Несколько процессов параллельно исследуют возможные продолжения порядка соединения, а продолжительность пробного выполнения ограничивается, чтобы медленный кандидат не задерживал дальнейший поиск. Здесь параллелизм относится к построению плана. На наборе JOB суммарное время выполнения запросов с LOOPLINE оказалось на 34,6\,\% меньше, чем с PostgreSQL~\cite{loopline2025}.
	
	В FOSS начальной точкой служит полный план, уже построенный штатным оптимизатором. Алгоритм последовательно проверяет два типа изменений: перестановку таблиц и замену способа выполнения конкретного соединения. Обучаемая модель оценивает преимущество одного варианта перед другим. На её основе авторы формируют моделируемую среду, в которой можно исследовать последовательности изменений без фактического выполнения каждого кандидата. На тестовых запросах JOB ускорение относительно PostgreSQL по суммарному времени составило 8,33 раза. На тестовой части TPC-DS оно составило 1,15 раза. TPC-DS --- стандартный тест аналитической обработки данных, моделирующий работу хранилища розничной сети. Он задаёт схему базы, правила генерации данных и 99 шаблонов SQL-запросов.\footnote{Transaction Processing Performance Council. \href{https://www.tpc.org/tpcds/}{Спецификация TPC-DS и описание теста}.} Таким образом, один и тот же способ изменения плана дал разный выигрыш на двух нагрузках~\cite{foss2024}.
	
	Следовательно, обучаемый компонент может вернуть прогноз числа строк, подсказку для построения плана или изменение уже выбранной операции. В первом случае качество оценивают по ошибке прогноза. В остальных случаях можно непосредственно измерить, за какое время СУБД выполнит запрос с полученным планом~\cite{wang2024,autoquo2024,loopline2025,foss2024}.
	
	
	\newpage
	\section{Оптимизация запросов в распределённых системах}
	
	В распределённой СУБД даже при известном порядке соединения нужно решить, где обработать строки с совпадающими ключами. Если эти строки находятся на разных узлах, часть данных потребуется передать по сети. Время операции будет зависеть от объёма пересылки, конкуренции за сетевые каналы и распределения работы между исполнителями~\cite{liang2024,coflop2023,prpq2017}.
	
	\subsection{Узлы выполнения, обмен и параллелизм}
	
	Первое решение возникает уже при чтении данных. Если фрагмент таблицы хранится на нескольких узлах, оптимизатор может выбрать, с какой копии его прочитать. Du и соавторы исследовали, как подобрать источники сразу для всех фрагментов запроса. Каждый вариант они представили в виде вектора, в котором для каждого фрагмента указан узел чтения. Для поиска авторы использовали модифицированный алгоритм искусственной пчелиной колонии DYABC-GO с операторами отбора, скрещивания и мутации. Качество варианта оценивалось по тому, насколько чтение нужных фрагментов удаётся сосредоточить на одних и тех же узлах. На шести синтетических запросах разной размерности DYABC-GO находил планы с меньшей расчётной стоимостью, чем сравниваемые методы MACGA и DEGA. При этом поиск занимал больше времени, чем у MACGA. Использованная функция стоимости напрямую не учитывала объём сетевого обмена и задержку передачи~\cite{du2024}.
	
	После выбора источников остаётся решить, какие строки передавать. Небольшую таблицу можно разослать всем исполнителям, чтобы каждый выполнил соединение на месте. Другой вариант --- перераспределить строки по ключу так, чтобы совпадающие значения попали к одному исполнителю. В первом случае объём обмена растёт с числом получателей копии, во втором зависит от того, сколько строк должно сменить узел~\cite{liang2024}.
	
	В RelJoin выбор между рассылкой и перераспределением зависит от относительной стоимости сетевой передачи и локальной обработки. Система использует доступную во время выполнения статистику о размерах обеих таблиц и учитывает степень параллелизма. Для сравнения авторы использовали механизм адаптивного выполнения запросов Spark (Adaptive Query Execution, AQE). На TPC-DS преимущество RelJoin было наибольшим при наименьшем из исследованных объёмов данных. Ниже эти результаты сопоставляются с экспериментом, в котором степень параллелизма подбирали отдельно для каждого способа соединения~\cite{liang2024,turkmen2026}.
	
	DPJoin развивает эту постановку. Каждой стратегии соединения соответствует своя нагрузка на сеть и процессоры. Поэтому DPJoin подбирает степень параллелизма для каждой стратегии с учётом объёма данных и уровня репликации, а затем сравнивает полученные варианты. При испытаниях на TPC-DS результат зависел от масштаба: при наибольшем исследованном масштабе меньшее среднее время показал AQE. Это ограничивает применимость результата одного эксперимента к данным другого масштаба~\cite{turkmen2026}.
	
	\subsection{Выполнение соединений и конкуренция за сеть}
	
	После выбора плана время операции зависит от того, как реализованы передача и локальная обработка строк. Barthels и соавторы сравнили распределённое хеш-соединение и соединение сортировкой и слиянием на системах с низкой сетевой задержкой. Авторы варьировали буферизацию, порядок отправки и способ доступа к удалённой памяти. На 4096 ядрах вариант хеш-соединения с предварительным разбиением данных достиг пропускной способности 48,7~млрд входных строк в секунду. Эта величина характеризует пропускную способность одного оператора~\cite{barthels2017}.
	
	Когда несколько соединений выполняются одновременно, их передачи могут конкурировать за одни и те же сетевые ресурсы. В системе CoFlop передачи, необходимые одному оператору, рассматриваются как группа потоков. Оператор сможет перейти к следующему этапу только после завершения всех необходимых ему передач, поэтому задержка одного потока влияет на время завершения группы. CoFlop учитывает размещение данных при выборе узлов назначения и согласует обслуживание конкурирующих групп. В экспериментах на сетевой модели вариант с меньшим объёмом переданных данных не всегда обеспечивал меньшее время завершения обмена. Для прогноза задержки существенны также направления и очерёдность передач~\cite{coflop2023}.
	
	\subsection{Перекос данных и балансировка нагрузки}
	
	При перераспределении по хешу строки с одинаковым ключом направляются одному получателю. При частом повторении одного ключа нагрузка на соответствующий узел может оказаться существенно выше. Часть исполнителей может завершить работу раньше, однако время всей операции будет определяться наиболее загруженным узлом. Следовательно, при фиксированном объёме таблиц изменение частот ключей способно изменить время выполнения того же физического плана~\cite{prpq2017,flowjoin2026}.
	
	В алгоритме PRPQ решение об обработке частых ключей принимается по локальной статистике. Для обычных значений выполняется перераспределение. Для часто встречающегося значения узел может сохранить строки у себя и запросить данные, необходимые для поиска совпадений. Это позволяет избежать пересылки всех повторяющихся строк одному получателю. Предварительное построение общей гистограммы по кластеру для такого решения не требуется. Авторы сопоставляли PRPQ, в частности, с PRPD, который распределяет и дублирует часть данных для частых ключей. На кластере из 16 узлов PRPQ снижал объём обмена и время соединения в исследованных сценариях~\cite{prpq2017}.
	
	В методе Patch-Based Shuffling группы строк назначаются узлам с учётом исходного размещения данных и уже распределённой нагрузки. Такой выбор позволяет обработать часть строк на месте, а часть передать менее загруженному исполнителю. Авторы оценивали эффект перераспределения в нескольких сценариях с разной выраженностью перекоса. В их модели время соединения включало сетевую пересылку и обработку данных. При слабом перекосе обычное хеш-разбиение могло выполняться быстрее. Для интерпретации результата существенно, что рабочие узлы моделировались потоками одного сервера~\cite{kassela2025}.
	
	Flow-Join++ связывает оценку кардинальности с балансировкой нагрузки. Модель заранее оценивает, сколько совпадающих строк приходится на отдельные значения ключа. Сразу после чтения таблиц алгоритм выделяет строки с ожидаемой высокой нагрузкой, чтобы согласовать их обработку на нескольких соединениях. Авторы сравнили метод со стандартными стратегиями на нагрузке Phone Bill, моделирующей аналитику в телекоммуникационной компании, и на модифицированных аналитических запросах теста TPC-H с управляемым перекосом. Отдельный практический вопрос состоит в том, сохраняется ли выигрыш при сравнении со штатным Spark AQE с включённой обработкой перекоса~\cite{flowjoin2026}.
	
	
	\newpage
	\section{Сравнение методов и возможности адаптации}
	
	Рассмотренные методы оценивают число строк, выбирают план или меняют выполнение отдельного оператора. Поэтому при сравнении важно указать, какое решение проверяли авторы и что измеряли в эксперименте. Это позволяет связать полученные результаты с выбором способа соединения в распределённой среде.
	
	\subsection{От расчётной стоимости плана к фактическому времени запроса}
	
	MAQO и AutoQuo подбирают порядок соединения, но оцениваются по разным критериям. Для MAQO можно установить, достигла ли эвристика минимальной расчётной стоимости. AutoQuo выдаёт PostgreSQL подсказку, после чего измеряется суммарное время полученных планов. Сокращение этого времени на 29,73\,\% характеризует эффект подсказок на запросах к IMDb, а вопрос о достижении глобального минимума стоимости остаётся открытым~\cite{goncalves2014,autoquo2024}.
	
	LOOPLINE и FOSS авторы также сравнивали со штатным PostgreSQL. LOOPLINE ищет продолжение порядка соединения и ограничивает длительность пробного запуска. FOSS начинает с готового плана и в моделируемой среде оценивает его изменения. В обоих исследованиях измеряли суммарное время выбранных планов, но каждый результат получен в собственном эксперименте. Прямое сравнение потребовало бы одинаковых запросов, конфигурации PostgreSQL и условий обучения. Время, затраченное на пробные запуски и построение планов, при этом следует учитывать отдельно~\cite{loopline2025,foss2024}.
	
	Влияние нагрузки можно отделить от различий между реализациями на примере самой FOSS. На тестовой части JOB выигрыш оказался выше, чем на TPC-DS (табл.~\ref{tab:plan-methods}). В JOB используются запросы с несколькими соединениями к базе фильмов IMDb. В TPC-DS запрос может, например, соединить сведения о продажах, товарах и датах, а затем сгруппировать результат. Для экспериментов с объёмом данных используется коэффициент масштаба SF (\textit{scale factor}). Он задаёт номинальный размер исходных данных, а фактический размер файлов зависит в том числе от формата хранения и сжатия. Поэтому одинаковый SF сам по себе ещё не гарантирует одинаковых условий сравнения~\cite{leis2015,foss2024,turkmen2026}.
	
	В экспериментах с RelJoin и DPJoin авторы использовали TPC-DS, чтобы проверить, как выбор распределённого соединения зависит от объёма данных. Liang и соавторы сравнили RelJoin с AQE в Spark. При SF~=~1 из 629 решений о способе соединения два метода разошлись в 66 случаях. В 63 из них RelJoin выбрал перераспределение строк по ключу, тогда как AQE предпочёл рассылку таблицы. На всей нагрузке RelJoin сэкономил 419,9~с, что соответствует снижению среднего времени запроса на 20,8\,\%. С ростом объёма данных преимущество сократилось до 2,0\,\% при SF~=~10 и до 1,3\,\% при SF~=~100. Эти результаты получены при отключённых в AQE объединении малых разделов и адаптивной обработке перекоса~\cite{liang2024}.
	
	В исследовании DPJoin три подхода сравнили в одном эксперименте. DPJoin подбирал для каждого способа соединения свою степень параллелизма. При SF~=~10 и коэффициенте репликации 3 это дало наименьшее среднее время на 99 запросах TPC-DS: 3,84~с у DPJoin, 4,09~с у AQE и 7,45~с у RelJoin. При SF~=~50 и том же коэффициенте репликации быстрее оказался AQE: 8,72~с против 10,25~с у DPJoin. Авторы связывают изменение результатов, в частности, с ростом нагрузки на память и затрат на сборку мусора. Этот пример показывает, почему даже при одном наборе запросов для сравнения методов важно указывать масштаб данных~\cite{turkmen2026}.
	
	В таблице~\ref{tab:plan-methods} наряду с результатами для полных запросов приведены показатели отдельных компонентов. Для RvCE измеряется ошибка прогноза числа строк, для реализаций Barthels и соавторов --- число входных строк, обработанных оператором соединения за секунду, для CoFlop --- время завершения группы сетевых передач. Связь с полным временем запроса в каждом случае различна. Например, более точный прогноз может помочь оптимизатору выбрать другой план, а ускорение сетевого обмена сократит время запроса, если дальнейшее выполнение ожидает завершения передачи~\cite{wang2024,barthels2017,coflop2023}.
	
	\begingroup
	\small
	\setlength{\tabcolsep}{4pt}
	\renewcommand{\arraystretch}{1.15}
	\begin{longtable}{@{}>{\raggedright\arraybackslash}p{.20\linewidth}>{\raggedright\arraybackslash}p{.34\linewidth}>{\raggedright\arraybackslash}p{\dimexpr.46\linewidth-4\tabcolsep\relax}@{}}
		\caption{Сравнение методов по задаче, условиям эксперимента и показателю результата}\label{tab:plan-methods}\\
		\toprule
		\textbf{Метод} & \textbf{Предмет исследования} & \textbf{Показатель и результат}\\
		\midrule
		\endfirsthead
		\multicolumn{3}{l}{Продолжение таблицы~\ref{tab:plan-methods}}\\
		\toprule
		\textbf{Метод} & \textbf{Предмет исследования} & \textbf{Показатель и результат}\\
		\midrule
		\endhead
		\bottomrule
		\endfoot
		MAQO~\cite{goncalves2014} & Поиск порядка соединения в H2. Авторы оценивают стоимость плана и время его выполнения. & В одном из экспериментов план MAQO выполнился почти на 23\,\% быстрее плана штатного H2. \\[2pt]
		AutoQuo~\cite{autoquo2024} & Частичный порядок как подсказка PostgreSQL. База --- штатный PostgreSQL, запросы к IMDb. & Суммарное время выполнения сократилось на 29,73\,\%. \\[2pt]
		LOOPLINE~\cite{loopline2025} & Параллельный поиск порядка соединения. База --- штатный PostgreSQL, набор JOB. & Суммарное время выполнения сократилось на 34,6\,\%. \\[2pt]
		FOSS~\cite{foss2024} & Изменение готового плана. База --- штатный PostgreSQL на каждом наборе запросов. & Ускорение на тестовой части по суммарному времени: 8,33 раза на JOB и 1,15 раза на TPC-DS. \\[2pt]
		RelJoin~\cite{liang2024} & Выбор способа соединения. База --- Spark AQE без объединения малых разделов и адаптивной обработки перекоса. & Среднее время на TPC-DS сократилось на 20,8\,\% при SF~=~1, на 2,0\,\% при SF~=~10 и на 1,3\,\% при SF~=~100. \\[2pt]
		DPJoin~\cite{turkmen2026} & Способ соединения и параллелизм. Прямое сравнение на TPC-DS при коэффициенте репликации 3. & Среднее время при SF~=~10: DPJoin --- 3,84~с, AQE --- 4,09~с, RelJoin --- 7,45~с. При SF~=~50: AQE --- 8,72~с, DPJoin --- 10,25~с. \\[2pt]
		RvCE~\cite{wang2024} & Прогноз числа строк сопоставляют с фактической кардинальностью. & На IMDb прогноз точнее, чем у сравниваемых обучаемых методов. Показатель --- $q$-error (1 при точном прогнозе). \\[2pt]
		Barthels и соавт.~\cite{barthels2017} & Распределённое хеш-соединение и соединение сортировкой и слиянием. & Пропускная способность хеш-соединения --- 48,7~млрд входных строк в секунду на 4096 ядрах. \\[2pt]
		CoFlop~\cite{coflop2023} & Планирование конкурирующих групп передач в сетевой модели. & Измеряют время до завершения последней передачи в группе. Меньший объём обмена не всегда давал меньшее время. \\[2pt]
	\end{longtable}
	\endgroup
	
	\subsection{Сопоставление методов обработки перекоса данных}
	
	При одинаковом масштабе данных время соединения зависит также от частот ключей. Обычное хеш-разбиение направляет все строки с одним ключом одному исполнителю. Если ключ встречается часто, завершение его задачи может задержать весь запрос. PRPQ, Patch-Based Shuffling и Flow-Join++ различаются тем, когда и на основании каких сведений они перераспределяют эту работу~\cite{prpq2017,kassela2025,flowjoin2026}.
	
	PRPQ использует локальные частоты ключей, тогда как для Patch-Based Shuffling важно также исходное размещение строк. Во втором сценарии Kassela и соавторов (рис.~15b исходной статьи) при параметре перекоса $\theta=1$ авторы оценили преимущество своего метода перед хеш-разбиением в 81\,\%. При минимальном перекосе в том же сценарии метод уступал хеш-разбиению. Этот результат нельзя напрямую сопоставить с PRPQ: для первого метода оценивали время в модели, для второго измеряли обмен и время соединения на физическом кластере~\cite{kassela2025,prpq2017}.
	
	Для цепочки соединений дополнительное значение имеет момент выявления частых ключей. Flow-Join++ оценивает число совпадений для ключей и использует прогноз при обработке нескольких соединений. Если дорогие ключи возникают лишь на одном этапе, выигрыш от раннего выделения может оказаться иным. Чтобы выяснить это, требуется сравнение на запросах одинаковой длины с разным расположением перекоса в плане. В таблице~\ref{tab:skew-methods} приведены границы экспериментальной проверки каждого подхода~\cite{prpq2017,kassela2025,flowjoin2026}.
	
	\begingroup
	\small
	\setlength{\tabcolsep}{4pt}
	\renewcommand{\arraystretch}{1.15}
	\begin{longtable}{@{}>{\raggedright\arraybackslash}p{.20\linewidth}>{\raggedright\arraybackslash}p{.34\linewidth}>{\raggedright\arraybackslash}p{\dimexpr.46\linewidth-4\tabcolsep\relax}@{}}
		\caption{Способы обработки перекоса данных и результаты экспериментов}\label{tab:skew-methods}\\
		\toprule
		\textbf{Метод} & \textbf{Принцип распределения работы} & \textbf{Результат и условия сравнения}\\
		\midrule
		\endfirsthead
		\multicolumn{3}{l}{Продолжение таблицы~\ref{tab:skew-methods}}\\
		\toprule
		\textbf{Метод} & \textbf{Как распределяется работа} & \textbf{Результат и условия сравнения}\\
		\midrule
		\endhead
		\bottomrule
		\endfoot
		PRPQ~\cite{prpq2017} & По локальным частотам узел выбирает, какие строки оставить у себя и для каких ключей запросить совпадения. & На физическом кластере из 16 узлов (192 ядра, до 2~млрд строк) метод сократил обмен и время соединения относительно исследованных альтернатив. \\[2pt]
		Patch-Based Shuffling~\cite{kassela2025} & Группы строк назначают с учётом размещения данных и загрузки исполнителей. & В сценарии~2 при $\theta=1$ преимущество по оценке времени составило 81\,\%. При минимальном перекосе выиграло хеш-разбиение. Узлы моделировали потоками одного сервера. \\[2pt]
		Flow-Join++~\cite{flowjoin2026} & Модель прогнозирует число совпадений по ключам. Тяжёлые ключи выделяются сразу после чтения для обработки в нескольких соединениях. & На Phone Bill и модифицированном TPC-H метод ускорял обработку относительно испытанных стратегий Spark.  \\[2pt]
	\end{longtable}
	\endgroup
	
	\subsection{Условия адаптации к распределённой обработке}
	
	Для распределённой системы можно развить принцип FOSS. Штатный оптимизатор строит исходный план. Затем порядок соединения таблиц фиксируют и сравнивают три способа обмена: рассылка малой таблицы, перераспределение по ключу или отдельная обработка частых значений. Такой эксперимент позволил бы применить идею локального изменения плана к выбору, который исследовали RelJoin и DPJoin. Здесь это предложение для дальнейшей работы~\cite{foss2024,liang2024,turkmen2026}.
	
	Проверяемая гипотеза состоит в том, что при выраженном перекосе выбор стратегии улучшится, если оценивать время самой долгой задачи наряду с общим объёмом передачи. Два плана могут пересылать одинаковый объём данных и по-разному загружать отдельные узлы. При слабом перекосе выигрыш от перераспределения может оказаться меньше времени, потраченного на сбор статистики. Та же зависимость между разовыми затратами и последующей экономией проявилась в работе AutoCox~\cite{coflop2023,prpq2017,kassela2025,flowjoin2026,rahman2026}.
	
	Для проверки гипотезы можно сравнить на одном кластере три варианта: штатный Spark AQE с включённой обработкой перекоса, выбор соединения по размерам таблиц и стоимости сети, а также тот же выбор с учётом частот ключей. Во втором и третьем вариантах следует сохранить одинаковыми порядок таблиц, доступные способы соединения и ресурсы. За основу можно взять запросы TPC-DS и варьировать распределение ключей при фиксированном размере входных таблиц. Изменение частот может изменить и число строк результата: его следует фиксировать там, где это возможно, либо измерять отдельно. Иначе ускорение может объясняться уменьшением объёма работы. Основной показатель --- полное время от поступления запроса до получения результата, включая сбор статистики. Для объяснения различий нужны также объём пересылки и длительность самой долгой задачи. Повторение измерений при нескольких масштабах данных позволит установить, когда учёт перекоса окупается.
	
	\newpage
	\section*{Заключение}
	\addcontentsline{toc}{section}{Заключение}
	
	В обзоре сопоставлены методы, которые решают три связанные задачи оптимизации запросов: оценивают размер промежуточных результатов, определяют затраты на их обработку и выбирают план из допустимых вариантов. Исследования классических оптимизаторов показывают, что итоговое время запроса зависит как от точности оценок, так и от состава рассмотренных планов. Более полный поиск может выявить удачный порядок соединения даже при неточной оценке числа строк. Вместе с тем расчётная стоимость плана остаётся полезным критерием лишь в той мере, в какой модель учитывает существенные для данной среды операции, например характер обращений к оперативной памяти~\cite{leis2015,manegold2000}.
	
	Обучаемые методы различаются тем, какое решение они передают оптимизатору. Для RvCE результатом служит оценка числа строк, для AutoQuo и LOOPLINE --- рекомендация порядка соединения, для FOSS --- изменение исходного плана. Поэтому даже точный прогноз кардинальности сам по себе не позволяет установить выигрыш во времени. В исследованиях, где измеряли запрос целиком, величина ускорения зависела от нагрузки. В частности, один и тот же метод FOSS дал существенно разный результат на тестовых запросах JOB и TPC-DS~\cite{wang2024,autoquo2024,loopline2025,foss2024}.
	
	В распределённой СУБД для одного порядка таблиц можно выбрать разные способы обмена и разную степень параллелизма. Эксперименты с RelJoin и DPJoin показывают, что при увеличении данных преимущество выбранной стратегии может уменьшиться или перейти к штатному оптимизатору. Работы о перекосе показывают, что при фиксированном объёме входных таблиц время соединения зависит и от того, сколько работы приходится на самого загруженного исполнителя. Эти результаты подводят к вопросу о том, при каких размерах таблиц и при какой степени перекоса выигрыш от отдельной обработки частых ключей покрывает затраты на сбор статистики и перераспределение строк~\cite{liang2024,turkmen2026,prpq2017,kassela2025,flowjoin2026}.
	
	Цель предполагаемого магистерского исследования --- разработать и экспериментально оценить метод выбора способа выполнения распределённого соединения с учётом частот ключей. Для достижения этой цели предполагается решить следующие задачи:
	\begin{enumerate} 
		\item Построить модель, которая по размерам таблиц и частотам ключей оценивает объём сетевого обмена, нагрузку наиболее загруженного узла и затраты на получение статистики. 
		\item Разработать на основе этой модели правило выбора между обычным хеш-перераспределением строк и отдельной обработкой часто встречающихся ключей. 
		\item Cравнить предложенный метод со штатной стратегией СУБД и с вариантом того же метода без учёта частот ключей, варьируя объём таблиц и степень перекоса.
\end{enumerate}
		 Основным показателем сравнения будет полное время выполнения запроса, включая сбор статистики. Такое сравнение позволит определить, в каких условиях сведения о частотах ключей меняют выбор стратегии и сокращают время выполнения запроса.
	
	% Разместить этот блок один раз в конце отчёта.
	% По мере написания следующих разделов добавлять новые \bibitem
	% в тот же список. Нумерация ссылок формируется автоматически.
	% Для обновления ссылок документ необходимо скомпилировать дважды.
	
	\providecommand{\refname}{}
	\renewcommand{\refname}{Список использованных источников}
	\providecommand{\bibname}{}
	\renewcommand{\bibname}{Список использованных источников}
	\newpage
	\begin{thebibliography}{99}
		\phantomsection\addcontentsline{toc}{section}{Список использованных источников}
		
		\bibitem{leis2015}
		Leis V., Gubichev A., Mirchev A., Boncz P., Kemper A., Neumann T.
		How Good Are Query Optimizers, Really? //
		Proceedings of the VLDB Endowment. ---
		2015. --- Vol.~9, no.~3. --- P.~204--215. ---
		DOI: \href{https://doi.org/10.14778/2850583.2850594}{10.14778/2850583.2850594}.
		
		\bibitem{manegold2000}
		Manegold S., Boncz P.~A., Kersten M.~L.
		Optimizing database architecture for the new bottleneck: memory access //
		The VLDB Journal. ---
		2000. --- Vol.~9, no.~3. --- P.~231--246. ---
		DOI: \href{https://doi.org/10.1007/s007780000031}{10.1007/s007780000031}.
		
		\bibitem{goncalves2014}
		Gonçalves F.~A.~C.~A., Guimarães F.~G., Souza M.~J.~F.
		Query join ordering optimization with evolutionary multi-agent systems //
		Expert Systems with Applications. ---
		2014. --- Vol.~41, no.~15. --- P.~6934--6944. ---
		DOI: \href{https://doi.org/10.1016/j.eswa.2014.05.005}{10.1016/j.eswa.2014.05.005}.
		
		\bibitem{wang2024}
		Wang Z., Duan H., Cheng Y., Min G.
		Learning complex predicates for cardinality estimation
		using recursive neural networks //
		Information Systems. ---
		2024. --- Vol.~124. --- Art.~102402. ---
		DOI: \href{https://doi.org/10.1016/j.is.2024.102402}{10.1016/j.is.2024.102402}.
		
		\bibitem{autoquo2024}
		Xiong X., Yu J., He Z.
		AutoQuo: An Adaptive plan optimizer with reinforcement learning for query plan selection //
		Knowledge-Based Systems. ---
		2024. --- Vol.~306. --- Art.~112664. ---
		DOI: \href{https://doi.org/10.1016/j.knosys.2024.112664}{10.1016/j.knosys.2024.112664}.
		
		\bibitem{foss2024}
		Zhong K., Sun L., Ji T., Li C., Chen H.
		FOSS: A Self-Learned Doctor for Query Optimizer //
		2024 IEEE 40th International Conference on Data Engineering (ICDE). ---
		2024. --- P.~4329--4342. ---
		DOI: \href{https://doi.org/10.1109/ICDE60146.2024.00330}{10.1109/ICDE60146.2024.00330}.
		
		\bibitem{liang2024}
		Liang F., Lau F.~C.~M., Cui H., Li Y., Lin B., Li C., Hu X.
		RelJoin: Relative-cost-based selection of distributed join methods
		for query plan optimization //
		Information Sciences. ---
		2024. --- Vol.~658. --- Art.~120022. ---
		DOI: \href{https://doi.org/10.1016/j.ins.2023.120022}{10.1016/j.ins.2023.120022}.
		
		\bibitem{prpq2017}
		Cheng L., Kotoulas S., Ward T.~E., Theodoropoulos G.
		Improving the robustness and performance of parallel joins over distributed systems //
		Journal of Parallel and Distributed Computing. ---
		2017. --- Vol.~109. --- P.~310--323. ---
		DOI: \href{https://doi.org/10.1016/j.jpdc.2017.06.016}{10.1016/j.jpdc.2017.06.016}.
		
		\bibitem{flowjoin2026}
		Panduman Y.~Y.~F., Ito R., Onizuka M.
		Flow-Join++: Skew-aware Distributed Joins with Cardinality-based
		Heavy Hitter Prediction //
		Journal of Information Processing. ---
		2026. --- Vol.~34. --- P.~671--685. ---
		DOI: \href{https://doi.org/10.2197/ipsjjip.34.671}{10.2197/ipsjjip.34.671}.
		
		\bibitem{rahman2026}
		Rahman M.~A., Hashmi S.~J., Gnanajothy K., Lee Y.-K.
		When temporary results meet intermediate index: An optimization technique
		of procedural SQL query processing //
		Data \& Knowledge Engineering. ---
		2026. --- Vol.~164. --- Art.~102577. ---
		DOI: \href{https://doi.org/10.1016/j.datak.2026.102577}{10.1016/j.datak.2026.102577}.
		
		\bibitem{loopline2025}
		Xiong X., Yu J., Han D., Tong X.
		LOOPLINE: Parallel Monte-Carlo tree search with knowledge constraints and its application in query optimization //
		Expert Systems with Applications. ---
		2025. --- Vol.~280. --- Art.~127490. ---
		DOI: \href{https://doi.org/10.1016/j.eswa.2025.127490}{10.1016/j.eswa.2025.127490}.
		
		\bibitem{coflop2023}
		Cheng L., Wang Y., Jhaveri R.~H., Wang Q., Mao Y.
		Toward Network-Aware Query Execution Systems in Large Datacenters //
		IEEE Transactions on Network and Service Management. ---
		2023. --- Vol.~20, no.~4. --- P.~4494--4504. ---
		DOI: \href{https://doi.org/10.1109/TNSM.2023.3273166}{10.1109/TNSM.2023.3273166}.
		
		\bibitem{du2024}
		Du Y., Cai Z., Ding Z.
		Query Optimization in Distributed Database Based on Improved Artificial Bee Colony Algorithm //
		Applied Sciences. ---
		2024. --- Vol.~14, no.~2. --- Art.~846. ---
		DOI: \href{https://doi.org/10.3390/app14020846}{10.3390/app14020846}.
		
		\bibitem{turkmen2026}
		Türkmen F., Ergenç Bostanoğlu B.
		Dynamic parallelization in distributed join optimization //
		PeerJ Computer Science. ---
		2026. --- Vol.~12. --- Art.~e3614. ---
		DOI: \href{https://doi.org/10.7717/peerj-cs.3614}{10.7717/peerj-cs.3614}.
		
		\bibitem{barthels2017}
		Barthels C., Müller I., Schneider T., Alonso G., Hoefler T.
		Distributed Join Algorithms on Thousands of Cores //
		Proceedings of the VLDB Endowment. ---
		2017. --- Vol.~10, no.~5. --- P.~517--528. ---
		DOI: \href{https://doi.org/10.14778/3055540.3055545}{10.14778/3055540.3055545}.
		
		\bibitem{kassela2025}
		Kassela E., Konstantinou I., Koziris N.
		Accelerating Distributed Repartition Joins on Skewed Datasets
		via Patch-Based Shuffling //
		IEEE Access. ---
		2025. --- Vol.~13. --- P.~41068--41096. ---
		DOI: \href{https://doi.org/10.1109/ACCESS.2025.3547623}{10.1109/ACCESS.2025.3547623}.
	\end{thebibliography}
	\includepdf[
	pages=-,
	fitpaper=true,
	pagecommand={\thispagestyle{empty}},
	pagecommand*={%
		\phantomsection
		\addcontentsline{toc}{section}{Приложение А. Сведения о публикациях}%
		\label{app:publications}%
	},
	picturecommand={%
		\put(0,32){\makebox[\paperwidth][c]{\normalfont\normalsize\thepage}}%
	}
	]{appendix_table.pdf}
\end{document}



